% corpus_supplement: Erdős problem #742
% source: https://www.erdosproblems.com/latex/742
% fetched_at: 2026-07-14T18:47:43Z
% payload_sha256: 76102e58688d13ac44595c327a7b500979b477438bd6d27439eaabf306be463f
% statement/remarks/references extracted by erdos_ingest.extract_source_page;
% rendered by erdos_ingest.render_tex — do not edit
\documentclass{article}
\usepackage{amsmath, amssymb, amsthm}
\usepackage{hyperref}
\usepackage{geometry}
\geometry{margin=1in}

\title{Erd\H{o}s Problem \#742}
\date{Last updated: 2025-08-31}
\author{Source: erdosproblems.com}

\begin{document}
\maketitle

\noindent\textbf{Status:} OPEN \quad \textbf{No prize}

\noindent\textbf{Tags:} graph theory

\medskip
\noindent\textbf{Problem Statement:}

Let $G$ be a graph on $n$ vertices with diameter $2$, such that deleting any edge increases the diameter of $G$. Is it true that $G$ has at most $n^2/4$ edges?

\bigskip
\noindent\textbf{Remarks:}

A conjecture of Murty and Plesnik (see \cite{CaHa79}) (although F\"{u}redi credits this conjecture to Murty and Simon, and further mentions that Erd\H{o}s told him that the conjecture goes back to Ore in the 1960s). The complete bipartite graph shows that this would be best possible.

This is true (for sufficiently large $n$) and was proved by F\"{u}redi \cite{Fu92}.

\bigskip
\noindent\textbf{References:}

[CaHa79] Caccetta, Louis and H\"{a}ggkvist, Roland, On diameter critical graphs. Discrete Math. (1979), 223-229.
 
 [Fu92] F\"{u}redi, Zolt\'{a}n, The maximum number of edges in a minimal graph of diameter
{$2$}. J. Graph Theory (1992), 81-98.

\bigskip
\noindent\small{Source: \url{https://www.erdosproblems.com/742}}
\end{document}
